\section{Ideal Esensial}
\begin{exam} Jika $A$ dan $B$ adalah aljabar-$C^*$, maka $A$ (lebih tepatnya, $A \oplus \{\vc
0\}$) merupakan ideal dalam $A \oplus B$.
\end{exam}

\begin{conv} Sebagaimana ditunjukkan oleh contoh sebelumnya, lazimnya $A$ dipandang sebagai
 \index{konvensi!$A$ merupakan subhimpunan dari $A \oplus B$}%
subhimpunan dari $A \oplus B$. Selanjutnya kita akan melakukan ini tanpa menyebutkannya lagi.
\end{conv}

\begin{notn} Untuk suatu elemen $c$ dari aljabar $A$
 \index{Ic@$I_c$}%
tetapkan
  \[ I_c := \biggl\{a_0c + cb_0 + \sum_{k=1}^p a_kcb_k \colon
   \text{$p \in \N$, $a_0$, \dots, $a_p$, $b_0$, \dots, $b_p \in A$}\biggr\} \]
\end{notn}

\begin{prop} Jika $c$ merupakan elemen dari aljabar $A$, maka $I_c$ merupakan ideal
(aljabar) dalam~$A$.
\end{prop}

Perhatikan bahwa dalam proposisi sebelumnya tidak dinyatakan bahwa ideal aljabar $I_c$
harus sejati. Sangat mungkin $I_c = A$ (misalnya, ketika $c$ merupakan
elemen invertibel dari aljabar beridentitas).

\vskip 2 pt

\begin{defn} Misalkan $c$ suatu elemen dari aljabar-$C^*$ $A$.
 \index{Jc@$J_c$ (ideal utama (tertutup) yang memuat~$c$)}%
 \index{utama!ideal}%
 \index{ideal!utama}%
Definisikan $J_c$, yakni \df{ideal utama} yang memuat $c$, sebagai irisan dari
keluarga semua ideal (tertutup) dalam $A$ yang memuat~$c$. Jelas bahwa $J_c$ merupakan ideal terkecil
yang memuat~$c$.
\end{defn}

\begin{prop} Dalam suatu aljabar-$C^*$, tutupan suatu ideal aljabar merupakan ideal.
\end{prop}

\begin{exam} Tutupan suatu ideal aljabar sejati dalam aljabar-$C^*$ tidak harus merupakan
ideal sejati. Sebagai contoh, $l_c$, yaitu himpunan barisan bilangan kompleks yang
pada akhirnya nol, padat dalam aljabar-$C^*$ $l_0 = \fml C_0(\N)$. (Namun, ingat
proposisi~\ref{0006801}.)
\end{exam}

\begin{prop} Jika $c$ suatu elemen dari aljabar-$C^*$, maka $J_c = \clo{I_c}$.
\end{prop}

\begin{notn} Kita menggunakan suatu konvensi notasi standar. Jika $A$ dan $B$ merupakan subhimpunan tak kosong
dari suatu aljabar, maka $AB$ berarti rentang linear dari hasil kali elemen-elemen dalam $A$ dan
elemen-elemen dalam $B$; yaitu,
 \index{<AB@$AB$ (dalam aljabar, rentang hasil kali elemen-elemen)}%
 \index{konvensi!dalam suatu aljabar, $AB$ menyatakan rentang hasil kali}%
$AB = \spn\{ab \colon \text{$a \in A$ dan $b \in B$}\}$. (Perhatikan bahwa dalam
definisi~\ref{000551} tidak ada bedanya apakah $AJ$ diartikan sebagai himpunan
hasil kali elemen-elemen dalam $A$ dengan elemen-elemen dalam $J$ atau rentang dari himpunan tersebut.)
\end{notn}

\begin{prop} Jika $I$ dan $J$ merupakan ideal dalam suatu aljabar-$C^*$, maka $\clo{IJ} =
I \cap J$.
\end{prop}

\vskip 2 pt

Suatu aljabar-$C^*$ tak beridentitas $A$ dapat dibenamkan sebagai ideal dalam aljabar-$C^*$ beridentitas melalui
berbagai cara. Aljabar-$C^*$ beridentitas terkecil yang memuat $A$ adalah unitalisasinya~$\wt
A$ (lihat proposisi~\ref{0015213}). Tentu saja tidak ada aljabar-$C^*$ beridentitas terbesar tempat
$A$ dapat dibenamkan sebagai ideal, sebab jika $A$ dibenamkan sebagai ideal dalam suatu
aljabar-$C^*$ beridentitas $B$ dan $C$ adalah \emph{sebarang} aljabar-$C^*$ beridentitas, maka $A$ merupakan ideal dalam
aljabar-$C^*$ beridentitas $B \oplus C$ yang lebih besar lagi. Alasan unitalisasi yang lebih besar ini
tidak begitu menarik adalah karena irisan ideal $C$ dengan $A$ ialah~$\{\vc 0\}$.
Hal ini memotivasi definisi berikut.

\vskip 2 pt

\begin{defn} Suatu ideal $J$ dalam aljabar-$C^*$ $A$ disebut
 \index{esensial!ideal}%
 \index{ideal!esensial}%
\df{esensial} jika dan hanya jika $I \cap J \ne \vc 0$ untuk setiap ideal tak nol $I$ dalam~$A$.
\end{defn}

\begin{exam} Suatu aljabar-$C^*$ $A$ merupakan ideal esensial dalam unitalisasinya $\wt A$ jika dan
hanya jika $A$ \emph{tidak} beridentitas.
\end{exam}

\begin{exam} Jika $H$ suatu ruang Hilbert, maka ideal operator kompak $\ofml K(H)$
merupakan ideal esensial dalam aljabar-$C^*$~$\ofml B(H)$.
\end{exam}

\begin{defn} Misalkan $J$ suatu ideal dalam aljabar-$C^*$~$A$. Kita mendefinisikan
$J^\perp$, yakni
 \index{anihilator}%
 \index{<superscript@$J^\perp$ (anihilator dari $J$)}%
\df{anihilator} dari $J$, sebagai $\bigl\{a \in A\colon Ja = \{\vc 0\}\bigr\}$.
\end{defn}

\begin{prop} Jika $J$ suatu ideal dalam aljabar-$C^*$, maka $J^\perp$ juga demikian.
\end{prop}

\begin{prop}\label{0038465} Suatu ideal $J$ dalam aljabar-$C^*$ $A$ bersifat esensial jika dan hanya
jika $J^\perp = \{\vc 0\}$.
\end{prop}

\begin{prop}\label{0038468} Jika $J$ suatu ideal dalam aljabar-$C^*$, maka $\bigl(J \oplus
J^\perp\bigr)^\perp = \{\vc 0\}$.
\end{prop}

\begin{notn} Misalkan $f$ suatu fungsi bernilai kompleks (atau real) pada himpunan $S$. Maka
  \[ Z_f := \{ s \in S\colon f(s) = 0\}. \]
Ini adalah
 \index{himpunan nol}%
 \index{Zf@$Z_f$ (himpunan nol dari~$f$)}%
\df{himpunan nol} dari~$f$.
\end{notn}

 \vskip 2 pt

Misalkan $A$ suatu aljabar-$C^*$ komutatif tak beridentitas. Menurut
\emph{teorema Gelfand--Naimark kedua}~\ref{00152156}, terdapat suatu ruang Hausdorff kompak lokal tak kompak
$X$ sedemikian sehingga $A = \fml C_0(X)$. (Di sini, tentu saja, kita membolehkan
diri melakukan penyalahgunaan bahasa yang lazim: secara harfiah, persamaan yang ditunjukkan itu
seharusnya berupa isomorfisme-$*\,$ isometrik.) Sekarang misalkan $Y$ suatu ruang Hausdorff kompak
yang memuat $X$ sebagai subhimpunan terbuka dan misalkan $B = \fml C(Y)$. Maka $B$ suatu aljabar-$C^*$
komutatif beridentitas. Pandang $A$ sebagai ideal dalam $B$ melalui pemetaan
$\iota\colon A \sto B\colon f \mapsto \wt f$ dengan
  \[ \wt f(y) = \left\{%
                \begin{array}{ll}
                         f(y), & \hbox{jika $y \in X$;} \\
                         0, & \hbox{selainnya.} \\
                \end{array}%
                \right.\]
Perhatikan bahwa himpunan tertutup $X^c$ adalah $\bigcap\{Z_{\wt f}\,\colon f \in \fml C_0(X)\}$.

\vskip 2 pt

\begin{prop}\label{0038474} Gunakan notasi seperti pada paragraf sebelumnya. Maka
ideal $A$ bersifat esensial dalam $B$ jika dan hanya jika subhimpunan terbuka $X$ padat dalam~$Y$.
\end{prop}

Dengan demikian, sifat esensial suatu ideal dalam konteks unitalisasi aljabar-$C^*$
komutatif tak beridentitas berkorespondensi secara tepat dengan sifat suatu subruang terbuka
yang padat dalam konteks kompaktifikasi ruang Hausdorff kompak lokal
tak kompak.





















\section{Kompaktifikasi dan Unitalisasi}
Dalam definisi~\ref{00152131}, objek yang keberadaannya dibuktikan dalam
proposisi sebelumnya~\ref{0015213} kita sebut ``unitalisasi'' dari suatu aljabar-$C^*$. Penyebutan seolah-olah tunggal
tersebut jelas menyesatkan. Sebagaimana ruang topologis dapat memiliki banyak
kompaktifikasi berbeda, suatu aljabar-$C^*$ dapat memiliki banyak unitalisasi. Jika $A$ merupakan
aljabar-$C^*$ komutatif tak beridentitas, jelas dari akibat~\ref{001924} bahwa
unitalisasi $\wt A$ adalah unitalisasi terkecil yang mungkin dari~$A$. Demikian pula, dalam topologi,
jika $X$ suatu ruang Hausdorff kompak lokal tak kompak, maka kompaktifikasi
satu titiknya jelas merupakan kompaktifikasi terkecil yang mungkin dari~$X$. Ingat bahwa
dalam proposisi~\ref{00152194} kita membuktikan bahwa membangun unitalisasi terkecil
dari $A$ ``pada dasarnya'' sama dengan membangun kompaktifikasi terkecil
dari~$X$.

 \vskip 2 pt

Kadang-kadang lebih nyaman (seperti, misalnya, pada paragraf sebelumnya) untuk mengambil suatu
``unitalisasi'' dari aljabar yang sudah beridentitas dan kadang-kadang lebih nyaman untuk mengambil suatu
``kompaktifikasi'' dari ruang yang sudah kompak. Karena tampaknya tidak ada
istilah yang diterima secara universal, saya memperkenalkan bahasa berikut (jelas tidak standar, tetapi
saya harap berguna).

\begin{defn}\label{0038614} Misalkan $A$ dan $B$ aljabar-$C^*$ serta $X$ dan $Y$
ruang topologis Hausdorff. Kita mengatakan bahwa
 \begin{enumerate}
    \item $B$ merupakan suatu \df{unitalisasi} dari $A$ jika $B$ beridentitas dan $A$ ($*\,$-isomorfik
dengan) suatu subaljabar-$C^*$ dari~$B$;
 \index{unitalisasi}%
 \index{unitalisasi!esensial}%
 \index{esensial!unitalisasi}%
    \item $B$ merupakan suatu \df{unitalisasi esensial} dari $A$ jika $B$ beridentitas dan $A$
($*\,$-isomorfik dengan) suatu ideal esensial dari~$B$;
 \index{kompaktifikasi}%
 \index{kompaktifikasi!esensial}%
 \index{esensial!kompaktifikasi}%
    \item $Y$ merupakan suatu \df{kompaktifikasi} dari $X$ jika ruang itu kompak dan $X$ (homeomorfik
dengan) suatu subruang dari~$Y$; dan
    \item $Y$ merupakan suatu \df{kompaktifikasi esensial} dari $X$ jika ruang itu kompak dan $X$
(homeomorfik dengan) suatu subruang padat dari~$Y$.
 \end{enumerate}
\end{defn}

Kiranya beberapa kata perlu disampaikan mengenai bagian-bagian dalam tanda kurung pada definisi
sebelumnya. Jarang terjadi bahwa suatu ruang topologis $X$ secara harfiah merupakan
\emph{subhimpunan} dari kompaktifikasi tertentu atas $X$ atau bahwa suatu aljabar-$C^*$ $A$ merupakan
\emph{subhimpunan} dari unitalisasi tertentu atas~$A$. Walaupun tentu benar bahwa sering kali
lebih nyaman untuk memandang satu aljabar-$C^*$ sebagai subhimpunan dari aljabar lain ketika sebenarnya yang
pertama hanya $*\,$-isomorfik dengan subhimpunan dari yang kedua, ada pula keadaan ketika
perincian menjadi lebih jelas jika pembenaman yang sesungguhnya ditentukan. Jika rincian perbedaan
ini belum sepenuhnya akrab, dua definisi berikut dimaksudkan untuk membantu.

\begin{defn} Misalkan $A$ dan $B$ aljabar-$C^*$. Kita mengatakan bahwa $A$
\df{dibenamkan} dalam~$B$ jika terdapat suatu homomorfisme-$*\,$ injektif $\iota\colon A \sto
B$; yaitu, jika $A$ isomorfik-$*\,$ dengan suatu subaljabar-$C^*$ dari~$B$ (lihat
proposisi~\ref{0019023} dan~\ref{00152181}). Homomorfisme-$*\,$ injektif $\iota$
merupakan suatu
 \index{C*@pembenaman-$C^*$}%
 \index{pembenaman!aljabar-$C^*$}%
\df{pembenaman} dari $A$ ke dalam~$B$. Dalam situasi ini lazimnya $A$ dan
jangkauan~$\iota$ diperlakukan sebagai aljabar-$C^*$ yang identik. Pasangan $(B,\iota)$ merupakan suatu
 \index{unitalisasi}%
\df{unitalisasi} dari~$A$ jika $B$ merupakan aljabar-$C^*$ beridentitas dan $\iota\colon A \sto B$ merupakan suatu
pembenaman. Unitalisasi $(B,\iota)$ bersifat
  \index{unitalisasi!esensial}%
  \index{esensial!unitalisasi}%
\df{esensial} jika jangkauan $\iota$ merupakan ideal esensial dalam~$B$.
\end{defn}

\begin{defn} Misalkan $X$ dan $Y$ ruang topologis Hausdorff. Kita mengatakan bahwa $X$
\df{dibenamkan} dalam~$Y$ jika terdapat homeomorfisme $j$ dari $X$ ke suatu subruang dari~$Y$.
Homeomorfisme $j$ merupakan suatu
 \index{topologis!pembenaman}%
 \index{pembenaman!topologis}%
\df{pembenaman} dari $X$ ke dalam~$Y$. Seperti dalam aljabar-$C^*$, lazimnya
jangkauan $j$ diidentifikasi dengan ruang~$X$. Pasangan $(Y,j)$ merupakan suatu
 \index{kompaktifikasi}%
\df{kompaktifikasi} dari~$X$ jika $Y$ suatu ruang Hausdorff kompak dan $j\colon X \sto Y$
merupakan suatu pembenaman. Kompaktifikasi $(Y,j)$ bersifat
  \index{kompaktifikasi!esensial}%
  \index{esensial!kompaktifikasi}%
\df{esensial} jika jangkauan $j$ padat dalam~$Y$.
\end{defn}

Kita telah membahas unitalisasi terkecil dari suatu aljabar-$C^*$ dan kompaktifikasi terkecil
dari suatu ruang Hausdorff kompak lokal. Sekarang, bagaimana dengan aljabar beridentitas terbesar, atau bahkan
maksimal, yang memuat~$A$? Jelas bahwa objek seperti itu tidak ada, sebab jika $B$ suatu
aljabar beridentitas yang memuat~$A$, maka $B \oplus C$ juga demikian, dengan $C$ sebarang
aljabar-$C^*$ beridentitas. Demikian pula, tidak ada ruang kompak terbesar yang memuat~$X$: jika $Y$ suatu
ruang kompak yang memuat $X$, maka gabungan saling lepas topologis $Y \uplus K$ juga demikian,
dengan $K$ sebarang ruang kompak tak kosong. Namun, masuk akal untuk menanyakan apakah
terdapat suatu unitalisasi esensial maksimal dari aljabar-$C^*$ atau suatu kompaktifikasi esensial
maksimal dari ruang Hausdorff kompak lokal. Jawabannya adalah \emph{ya} dalam kedua
kasus. \emph{Kompaktifikasi Stone--\v{C}ech} $\beta(X)$ yang terkenal bersifat maksimal di antara
kompaktifikasi esensial dari ruang Hausdorff kompak lokal tak kompak~$X$. Perinciannya
dapat ditemukan dalam sebarang buku teks topologi yang baik. Salah satu uraian standar yang mudah dibaca
adalah~\cite{Willard:1968}, butir 19.3--19.12. Pendekatan yang lebih canggih menggunakan sebagian
analisis fungsional---lihat, misalnya, \cite{Conway:1990}, bab V, bagian~6. Ternyata
terdapat pula unitalisasi esensial maksimal dari suatu aljabar-$C^*$ tak beridentitas
$A$---objek ini disebut \emph{aljabar pengali} dari~$A$.

Kita mengatakan bahwa suatu unitalisasi esensial $M$ dari aljabar-$C^*$ $A$ bersifat \emph{maksimal} jika setiap
aljabar-$C^*$ yang memuat $A$ sebagai ideal esensial dapat dibenamkan dalam~$M$. Berikut pernyataan yang lebih
formal.

\begin{defn} Suatu unitalisasi esensial $(M,j)$ dari aljabar-$C^*$ $A$ disebut
 \index{maksimal!unitalisasi esensial}%
 \index{esensial!unitalisasi!maksimal}%
 \index{unitalisasi!esensial maksimal}%
\df{maksimal} jika untuk setiap pembenaman $\iota\colon A \sto B$ yang jangkauannya merupakan ideal esensial
dalam~$B$, terdapat homomorfisme-$*\,$ $\phi\colon B \sto M$ sedemikian sehingga $\phi
\circ \iota = j$.
\end{defn}

\begin{prop} Dalam definisi sebelumnya, homomorfisme-$*\,$ $\phi$, jika ada, harus
injektif.
\end{prop}

\begin{prop} Dalam definisi sebelumnya, homomorfisme-$*\,$ $\phi$, jika ada, harus
tunggal.
\end{prop}

Bandingkan definisi berikut dengan~\ref{0028}.

\begin{defn} Misalkan $A$ dan $B$ aljabar-$C^*$, dan $V$ modul Hilbert-$A$. Pemetaan
$\phi\colon B \sto \ofml L(V)$ yang merupakan homomorfisme-$*\,$ disebut
 \index{tak terdegenerasi!homomorfisme-$*\,$}%
\df{tak terdegenerasi} jika $\phi^\sto(B)\,V$ padat dalam~$V$.
\end{defn}

\begin{prop} Misalkan $A$, $B$, dan $J$ aljabar-$C^*$, $V$ suatu modul Hilbert-$B$, serta
$\iota\colon J \sto A$ suatu homomorfisme-$*\,$ injektif yang jangkauannya merupakan ideal dalam~$A$.
Jika $\phi\colon J \sto \ofml L(V)$ suatu homomorfisme-$*\,$ tak terdegenerasi, maka terdapat
perpanjangan tunggal dari $\phi$ menjadi homomorfisme-$*\,$ $\overline\phi\colon A \sto
\ofml L(V)$ yang memenuhi $\overline\phi \circ \iota = \phi$.
\end{prop}

\begin{prop}\label{0038621} Jika $A$ suatu aljabar-$C^*$ tak nol, maka $(\ofml L(A),L)$ merupakan
unitalisasi esensial maksimal dari~$A$. Unitalisasi ini tunggal dalam arti bahwa jika $(M,j)$ merupakan
unitalisasi esensial maksimal lain dari $A$, maka terdapat suatu isomorfisme-$*\,$ $\phi
\colon M \sto \ofml L(A)$ sedemikian sehingga $\phi \circ j = L$.
\end{prop}

\begin{defn}\label{0038641} Misalkan $A$ suatu aljabar-$C^*$. Kita mendefinisikan
 \index{aljabar pengali}%
 \index{aljabar!pengali}%
\df{aljabar pengali} dari $A$ sebagai keluarga $\ofml L(A)$ operator yang dapat diadjoinkan
pada~$A$. Mulai sekarang kita menyatakan keluarga ini
 \index{MA@$M(A)$ (aljabar pengali dari $A$)}%
dengan $M(A)$.
\end{defn}









\endinput
